% ============================================================
\section{行为效用、社会偏好与离散选择}
% ============================================================

\subsection{参照依赖价值函数（前景理论型）}

\notehead{形式与性质}
\[
v(z)=
\begin{cases}
z^{\alpha}, & z\ge0,\\[2pt]
-\lambda(-z)^{\beta}, & z<0,
\end{cases}
\qquad \alpha,\beta\in(0,1),\ \lambda>1.
\]
三个结构性特征，各对应一类实验事实：
\begin{enumerate}
  \item \textbf{参照依赖}：自变量是\textbf{相对参照点的损益} $z$，不是财富水平；
  \item \textbf{损失厌恶}：$\lambda>1$，损失的效用权重是等额收益的 $\lambda$ 倍
        （估计值常落在 $2$ 附近）；
  \item \textbf{敏感度递减}：$\alpha,\beta<1$，损益越大边际感受越小——
        价值函数在收益域凹、损失域凸（S 形）。
\end{enumerate}
直接推论：\textbf{反射效应}——收益域风险规避、损失域风险偏好
（“$+100$ 确定”对“$+200$ 半概率”，多数人选前者；
“$-100$ 确定”对“$-200$ 半概率”，多数人选后者）。
期望效用无法在一个 $u$ 下同时产生这两个选择。

\begin{center}
\begin{tikzpicture}[scale=0.95]
  \draw[axisline] (-3.4,0) -- (3.6,0) node[below left, axislabel] {$z$（相对参照点）};
  \draw[axisline] (0,-3.0) -- (0,2.6) node[below left, axislabel] {$v(z)$};
  \draw[thick, MidBlue, domain=0:3.2, smooth, variable=\x] plot ({\x},{sqrt(\x)*1.15});
  \draw[thick, MidBlue, domain=0:3.2, smooth, variable=\x] plot ({-\x},{-(\x)^0.6*1.55});
  \node[nodeblue] at (0,0) {};
  \draw[dashed, MidBlue] (0,0) -- (2.9,2.35) node[right, align=left] {损失域切线更陡：\\ 斜率比 $\approx\lambda$};
  \node[align=left] at (2.3,-1.5) {原点折点（损失厌恶）\\ 收益凹、损失凸\\ （敏感度递减）};
\end{tikzpicture}
\end{center}
\figexplan{横轴为相对参照点的损益 $z$，纵轴为价值 $v(z)$。曲线在原点出现折点：收益域为凹（风险厌恶）、损失域为凸（风险寻求），且损失域切线更陡（虚线），斜率之比约为损失厌恶系数 $\lambda$。原点处的折点即参照依赖，两侧曲率递减对应敏感度递减。}

\notehead{适用情形}
资产配置（处置效应：卖盈持亏）、劳动供给（日收入目标）、
定价与营销（免费试用后的退订损失）、保单设计。
与概率权重函数配合即为完整前景理论；仅价值函数已能解释大多数参照点现象。

\subsection{不平等厌恶（收益分配型社会偏好）}\label{subsec:fs}

\notehead{形式一：差异厌恶}
\[
u_i=x_i-\alpha_i\frac{1}{n-1}\sum_{j\ne i}\max\{x_j-x_i,\,0\}
-\beta_i\frac{1}{n-1}\sum_{j\ne i}\max\{x_i-x_j,\,0\},
\qquad 0\le\beta_i\le\alpha_i.
\]
$\alpha_i$：对\textbf{不如人}的厌恶（劣势不平等）；
$\beta_i$：对\textbf{优于人}的不安（优势不平等），且设定 $\beta_i\le\alpha_i$
——不如人的痛苦强于优于人的不安，这是实验数据的稳健结论。
两人情形简化为
\[
u_i=
\begin{cases}
x_i-\beta_i(x_i-x_j), & x_i\ge x_j,\\[2pt]
x_i-\alpha_i(x_j-x_i), & x_i<x_j.
\end{cases}
\]
无差异曲线在 $x_i=x_j$ 处成折线，折点两侧斜率不同（$\alpha$ 侧更陡）。

\begin{center}
\begin{tikzpicture}[scale=1.0]
  \draw[axisline] (0,0) -- (5.6,0) node[below left, axislabel] {$x_i$};
  \draw[axisline] (0,0) -- (0,4.4) node[below left, axislabel] {$x_j$};
  \draw[dashed, MidBlue] (0,0) -- (4.2,4.2) node[above right] {$x_i=x_j$};
  % kinked ICs: lines with kink on diagonal
  \draw[thick, MidBlue] (0.5,3.8) -- (2.2,2.2) -- (4.9,1.5);
  \draw[thick, MidBlue] (0.5,2.6) -- (1.4,1.4) -- (4.9,0.35);
  \node[align=left] at (4.85,3.05) {折线无差异曲线：\\ 对角线左侧（不如人）\\ 斜率更陡（$\alpha$ 侧）};
\end{tikzpicture}
\end{center}
\figexplan{虚线为 $x_i=x_j$ 的平等线。无差异曲线在该线上折曲：对角线左侧（自己不如他人）斜率更陡，对应不利不平等厌恶参数 $\alpha$；右侧斜率较缓，对应有利不平等厌恶 $\beta$。折点两侧斜率不等，正是分配偏好无法用光滑效用表示的表现。}

\notehead{形式二：社会最优与最小者关心}
\[
u_i=(1-\rho-\sigma)x_i+\rho\min(x_i,x_j)+\sigma\max(x_i,x_j).
\]
$\rho>0$ 时在意\textbf{最差者}（慈善/效率兼顾取决于 $\sigma$ 符号与大小）；
$\sigma<0$ 时甚至\textbf{嫉妒}（宁损己以损人）。
与差异厌恶形式的区别：本形式关心极端值而非逐对差距，
对分配中的“最底层”更敏感。

\notehead{适用情形}
最后通牒/独裁者/公共品博弈的结构估计、
组织内薪酬压缩的福利分析、再分配偏好的调查实验。
估计要点：$\alpha$ 由拒绝行为识别，$\beta$ 由分配行为识别。

\subsection{离散选择：随机效用模型}\label{subsec:rum}

\notehead{形式与 Logit 推导}
个体 $i$ 对方案 $j$ 的效用
\[
u_{ij}=v_j(\boldsymbol{x}_j;\theta)+\varepsilon_{ij},
\]
$v_j$ 为系统部分（可观测属性的函数），$\varepsilon_{ij}$ 为研究者不可观测的口味扰动。
设 $\{\varepsilon_{ij}\}$ 独立同分布于 I 型极值分布 $F(\varepsilon)=\exp(-e^{-\varepsilon})$，则
\[
\begin{aligned}
P_i(j)&=\Pr\bigl(v_j+\varepsilon_{ij}\ge v_k+\varepsilon_{ik},\ \forall k\bigr)\\
&=\int\prod_{k\ne j}F\bigl(v_j-v_k+\varepsilon\bigr)\,f(\varepsilon)\,\mathrm{d}\varepsilon
=\boxed{\ \frac{e^{v_j}}{\sum_{k}e^{v_k}}\ }.
\end{aligned}
\]
关键一步：I 型极值分布的稳定性——独立极值变量之差服从 Logistic 分布，
多重积分闭式可解。选择概率只依赖\textbf{效用差}，与绝对水平无关。

\notehead{性质与扩展}
\begin{itemize}
  \item \textbf{IIA 性质}：两方案的选择概率比 $P(j)/P(k)=e^{v_j-v_k}$ 与第三方无关——
        既是计算优点也是经验缺陷（红蓝巴士悖论）；
  \item \textbf{嵌套 Logit}：方案分组（巢内相关），放松 IIA——
        与 §\ref{subsec:flexible} 分支树型效用互为表里；
  \item \textbf{随机系数 Logit}：$\theta$ 在个体间随机分布，
        允许替代模式任意化，是需求估计的工业标准（§\ref{subsec:jpe-markups}）。
\end{itemize}

\notehead{适用情形}
商品选择（品牌、车型）、区位选择、职业与学校选择、
任何“多选一”且需要\textbf{从选择数据反推偏好}的经验场景。
福利计算有闭式公式：期望最大效用（消费者剩余）
$\mathbb{E}\max_j u_{ij}=\ln\sum_j e^{v_j}+$ 常数（对数和公式）。

\subsection{类型学小结：六个选择维度}

\begin{center}
\small
\begin{tabular}{llll}
\toprule
维度 & 问题 & 形式分支 & 代表 \\
\midrule
商品结构 & 替代还是互补？ & 线性 / 乘积型 / CES / 固定比例 & §4.1 \\
收入效应 & 份额恒定还是变动？ & 位似 / 基本需求型 / 非位似 & §4.1.5--§4.1.7, §4.2 \\
财富效应 & 是否可忽略？ & 拟线性 vs 一般 & §4.1.6 \\
时间 & 一致还是不一致？ & 指数贴现 / $\beta$-$\delta$ & §4.3.1, §4.3.4 \\
风险 & 时间可加还是解绑？ & CRRA / 递归型 / 失望厌恶 & §4.3.3, §4.3.5--§4.3.6 \\
他人 & 自利的还是社会的？ & 自利 / 差异厌恶 / 攀比 & §4.3.2, §4.5.2 \\
\bottomrule
\end{tabular}
\end{center}
